% ─────────────────────────────────────────────────────────────────────────────
\section{2D Benchmark Results: NAS Architecture Transfer and Adaptive-$k$}
\label{sec:results_2d_adaptive}

This section evaluates the adaptive-$k$ MSWP strategy on four 2D benchmark
domains.
It compares two architecture sources: the transfer architecture taken directly
from the NAS-PINN reference work~\cite{wang2023}, and the three architectures
found by our own NAS searches.
All experiments use the same physical parameters, benchmark geometries, and
training protocol.

% ─────────────────────────────────────────────────────────────────────────────
\subsection{Domains, Parameters, and Training Setup}
\label{subsec:2d_adap_setup}

We used four 2D benchmark domains.
The Square and Circle have regular boundaries.
The L-Shape has a 270° reentrant corner that concentrates the temperature
gradient locally.
The Flower has multiple concave boundary regions and is the most geometrically
complex 2D domain in this study.
Together they form a progression from the simplest to the most irregular
boundary.

The same physical and training parameters were used for all four domains and
all four architectures.
Table~\ref{tab:2d_adap_params} lists these parameters.

\begin{table}[H]
\centering
\caption{Physical and training parameters for the 2D adaptive-$k$ experiments.}
\label{tab:2d_adap_params}
\small
\begin{tabular}{llc}
\toprule
Parameter & Symbol & Value \\
\midrule
Thermal conductivity     & $k_T$      & $160\,\mathrm{W/(mK)}$ \\
Volumetric heat capacity & $\rho c_p$ & $2.4\times10^6\,\mathrm{J/(m^3K)}$ \\
Convection coefficient   & $h$        & $5000\,\mathrm{W/(m^2K)}$ \\
Initial temperature      & $T_0$      & $540\,^\circ\mathrm{C}$ \\
Quench bath temperature  & $T_w$      & $20\,^\circ\mathrm{C}$ \\
\midrule
Adam epochs (Bayesian/TPE)    & --- & 1\,000 \\
Adam epochs (NSGA-II/III)     & --- & 1\,000 \\
L-BFGS iterations per window  & --- & 300 \\
Interior collocation points   & $n_\mathrm{domain}$ & 1\,500 \\
Boundary collocation points   & $n_\mathrm{bc}$     & 300 \\
Adaptive-$k$ promotion threshold  & $\tau_\mathrm{up}$   & $2.0\,^\circ\mathrm{C}$ \\
Adaptive-$k$ demotion threshold   & $\tau_\mathrm{down}$ & $5.0\,^\circ\mathrm{C}$ \\
\bottomrule
\end{tabular}
\end{table}

The NAS transfer architecture $[90,70,30,110]$ was taken directly from Wang
et al.~\cite{wang2023} and tested at seed~0 only.
The three NAS-found architectures (Bayesian/TPE $[151{\times}5]$, NSGA-II
$[153{\times}3]$, NSGA-III $[75{\times}3]$) were each trained with ten
independent random seeds.

Table~\ref{tab:2d_adap_combined} gives the complete picture for all four
domains and all four architectures.
Each row shows mean-window MAE, $L_2$ relative error, FEM saving, training
time, and the adaptive $k$-schedule the controller selected.
The $k$-schedule column lists the skip factor used at each consecutive
prediction window from $t=0$ to $t=30\,\mathrm{s}$.
NAS Transfer results are from a single seed (seed~0) because this architecture
serves as a reference baseline only.
The three NAS-found architectures report ten-seed mean~$\pm$~std.
Full per-seed individual values are in Appendix~\ref{app:full_results}
(Tables~\ref{tab:adap_square_seeds}--\ref{tab:adap_flower_seeds}).

\begin{table}[H]
\centering
\footnotesize
\setlength{\tabcolsep}{3pt}
\caption{2D adaptive-$k$ experiments: all four domains and architectures.
MAE ($^\circ\mathrm{C}$): NAS Transfer = seed~0 only (single baseline run);
Bayesian/TPE, NSGA-II, NSGA-III = ten-seed mean~$\pm$~std.
Train (s): total training time (seed~0 for NAS Transfer; mean for others).
$k$-schedule: skip factor at each consecutive prediction window (seed~0).
\textbf{Bold} = lowest MAE per domain.}
\label{tab:2d_adap_combined}
\resizebox{\textwidth}{!}{%
\begin{tabular}{@{}llccccc@{}}
\toprule
Domain & Architecture & MAE ($^\circ\mathrm{C}$) & $L_2$ & FEM saved & Train (s) & $k$-schedule \\
\midrule
\multirow{4}{*}{Square}
  & NAS Transfer (seed~0) & 1.662              & 0.016 & 85\% & 471 & 1-1-1-1-2-3-4-5-5 \\
  & Bayesian/TPE          & $1.641\pm0.104$    & 0.017 & 85\% & 503 & 1-1-1-1-2-3-4-5-5 \\
  & NSGA-II               & $1.768\pm0.059$    & 0.015 & 85\% & 420 & 1-1-1-1-2-3-4-5-5 \\
  & NSGA-III              & $1.717\pm0.047$    & 0.015 & 85\% & 425 & 1-1-1-1-2-3-4-5-5 \\
\midrule
\multirow{4}{*}{Circle}
  & NAS Transfer (seed~0) & 1.477              & 0.014 & 85\% & 484 & 1-1-1-1-2-3-4-5-5 \\
  & Bayesian/TPE          & $1.527\pm0.105$    & 0.017 & 85\% & 546 & 1-1-1-1-2-3-4-5-5 \\
  & NSGA-II               & $1.648\pm0.035$    & 0.014 & 85\% & 469 & 1-1-1-1-2-3-4-5-5 \\
  & NSGA-III              & $1.619\pm0.029$    & 0.013 & 85\% & 465 & 1-1-1-1-2-3-4-5-5 \\
\midrule
\multirow{4}{*}{L-Shape}
  & NAS Transfer (seed~0) & 2.378              & 0.022 & 83\% & 518 & 1-1-1-1-1-1-2-3-4-5 \\
  & Bayesian/TPE          & $\mathbf{1.357\pm0.036}$ & 0.017 & 85\% & 506 & 1-1-1-1-2-2-3-4-5$^*$ \\
  & NSGA-II               & $2.438\pm0.052$    & 0.023 & 83\% & 491 & 1-1-1-1-1-2-3-4-5-5$^*$ \\
  & NSGA-III              & $2.344\pm0.032$    & 0.021 & 83\% & 475 & 1-1-1-1-1-1-2-3-4-5 \\
\midrule
\multirow{4}{*}{Flower}
  & NAS Transfer (seed~0) & 2.007              & 0.017 & 85\% & 679 & 1-1-1-1-2-3-4-5-5 \\
  & Bayesian/TPE          & $\mathbf{1.741\pm0.105}$ & 0.015 & 85\% & 722 & 1-1-1-1-2-3-4-5-5 \\
  & NSGA-II               & $2.113\pm0.054$    & 0.016 & 85\% & 669 & 1-1-1-1-2-3-4-5-5 \\
  & NSGA-III              & $2.001\pm0.055$    & 0.016 & 85\% & 653 & 1-1-1-1-2-3-4-5-5 \\
\bottomrule
\multicolumn{7}{l}{\small $^*$L-Shape schedules vary across seeds. See Appendix~\ref{app:full_results} (Table~\ref{tab:lshape_k_schedules}).}
\end{tabular}}
\end{table}

\begin{figure}[H]
\centering
\begin{subfigure}[t]{0.49\textwidth}
  \includegraphics[width=\linewidth]{fem_pinn_adaptive_k_square.png}
  \caption{Square}
  \label{fig:adap2d_square}
\end{subfigure}\hfill
\begin{subfigure}[t]{0.49\textwidth}
  \includegraphics[width=\linewidth]{fem_pinn_adaptive_k_circle.png}
  \caption{Circle}
  \label{fig:adap2d_circle}
\end{subfigure}\\[6pt]
\begin{subfigure}[t]{0.49\textwidth}
  \includegraphics[width=\linewidth]{fem_pinn_adaptive_k_lshape.png}
  \caption{L-Shape}
  \label{fig:adap2d_lshape}
\end{subfigure}\hfill
\begin{subfigure}[t]{0.49\textwidth}
  \includegraphics[width=\linewidth]{fem_pinn_adaptive_k_flower.png}
  \caption{Flower}
  \label{fig:adap2d_flower}
\end{subfigure}
\caption{FEM vs PINN mean temperature trajectories and per-window MAE for all
four 2D domains under adaptive-$k$ control (Bayesian/TPE, seed~0).
The shaded band shows $\pm$one standard deviation across seeds.}
\label{fig:adap2d_all}
\end{figure}

On the Square and Circle domains, all four architectures produce nearly
identical results.
MAE differences are below $0.3\,^\circ\mathrm{C}$ and all configurations
achieve 85\% FEM saving.
The NAS transfer architecture generalises well to these smooth geometries.

On the L-Shape domain, Bayesian/TPE achieves $1.357\pm0.036\,^\circ\mathrm{C}$
while the NAS transfer and both NSGA variants reach $2.34$--$2.46\,^\circ\mathrm{C}$.
The same pattern appears on the Flower domain.
The reentrant corner and concave boundaries concentrate temperature gradients
locally.
The larger capacity of Bayesian/TPE (${\approx}93$K parameters, ReLU)
better represents these gradients.
This shows that geometry-specific NAS adds value when the target domain
contains features absent from the NAS validation domain.

The ten-seed MAE standard deviation remains below $0.10\,^\circ\mathrm{C}$
on Square, Circle, and Flower.
This low variance confirms that training is stable and results are
reproducible.
On Square, Circle, and Flower, all ten seeds produce the identical
$k$-schedule, confirming that the adaptive controller responds to the
physics of the quench and not to random initialisation.
On the L-Shape, schedules vary by one window across seeds for Bayesian/TPE
and NSGA-II.
The concentrated gradient at the reentrant corner makes the first-window MAE
sensitive to weight initialisation, which shifts the promotion point by one
window in some seeds.
The full per-seed schedule breakdown is in Appendix~\ref{app:full_results}
(Table~\ref{tab:lshape_k_schedules}).

% ─────────────────────────────────────────────────────────────────────────────
\subsection{FEM vs PINN Temperature Comparison}
\label{subsec:2d_adap_fem_pinn}

Figure~\ref{fig:adap2d_all} shows the FEM vs PINN mean temperature trajectories
and per-window MAE for all four domains.
In all four domains, the PINN mean temperature trajectories closely follow
the FEM reference throughout the 30-second quench.
The MAE is highest in the first one to two windows where the quenching
gradient is steepest, then decreases as the field relaxes.
On the Square and Circle domains, all three architectures perform similarly
and remain well below $3\,^\circ\mathrm{C}$ after the initial transient.
On the L-Shape and Flower domains, Bayesian/TPE achieves lower MAE than
NSGA-II and NSGA-III, consistent with the scalar metrics in
Table~\ref{tab:2d_adap_combined}.

% ─────────────────────────────────────────────────────────────────────────────
\subsection{PINN-Only Rollout}
\label{subsec:2d_adap_pinn_only}

PINN-only rollout (without periodic FEM correction) was not evaluated
systematically on these four domains.
A preliminary experiment on the Square domain using a smaller architecture
produced a mean-window MAE of $29.6\,^\circ\mathrm{C}$, compared to
$1.64\,^\circ\mathrm{C}$ for the FEM-anchored adaptive-$k$ run.
This is an 18-fold increase.
The comprehensive PINN-only analysis on the Thermal Fin benchmark
(Section~\ref{subsec:tf_pinn_only}) confirms this failure pattern: without
FEM correction, errors accumulate rapidly and reach four to eighteen times
the engineering acceptance threshold.
These results establish that periodic FEM anchoring is essential for accurate
predictions on all benchmark domains, not only the Thermal Fin.

In summary, the 2D adaptive-$k$ results lead to three conclusions.
First, the NAS transfer architecture and our NAS-found architectures produce
comparable results on smooth 2D geometries.
Geometry-specific NAS adds value only when the target domain contains
features such as reentrant corners or concave boundaries that are absent
from the NAS validation domain.
Second, all three NAS-found architectures achieve 83--85\% FEM saving with
MAE below $2.5\,^\circ\mathrm{C}$ across all four domains.
The ten-seed standard deviation remains below $0.10\,^\circ\mathrm{C}$.
Third, FEM anchoring is essential.
Without periodic correction, errors accumulate to levels that are not
useful for engineering predictions.

The next section presents the systematic fixed-$k$ sweep on the three
canonical 2D domains, which provides a detailed accuracy--efficiency
trade-off analysis across all skip factors $k=1$ to $k=5$.
